%% LyX 2.4.4 created this file.  For more info, see https://www.lyx.org/.
%% Do not edit unless you really know what you are doing.
\documentclass[english]{article}
\usepackage[T1]{fontenc}
\usepackage[latin9]{inputenc}
\usepackage{enumitem}
\usepackage{amsmath}
\usepackage{amsthm}
\usepackage{amssymb}
\usepackage{geometry}
\geometry{verbose,tmargin=1in,bmargin=1in,lmargin=1in,rmargin=1in}
\PassOptionsToPackage{normalem}{ulem}
\usepackage{ulem}

\makeatletter
%%%%%%%%%%%%%%%%%%%%%%%%%%%%%% Textclass specific LaTeX commands.
\numberwithin{equation}{section}
\numberwithin{figure}{section}
\newlength{\lyxlabelwidth}      % auxiliary length 
\theoremstyle{plain}
\newtheorem{thm}{\protect\theoremname}[section]
\theoremstyle{definition}
\newtheorem{xca}[thm]{\protect\exercisename}

%%%%%%%%%%%%%%%%%%%%%%%%%%%%%% User specified LaTeX commands.
\usepackage{fancyhdr}
\usepackage{lastpage}
\pagestyle{fancy}
\usepackage{enumitem}
\usepackage{ifthen}
\everymath=\expandafter{\the\everymath\displaystyle}
%\DeclareMathSizes{10}{10}{10}{10}
\usepackage{multicol}

\usepackage{tikz}
\usetikzlibrary{patterns}
\usetikzlibrary{plotmarks}
\usepackage{pgfplots}
\pgfplotsset{compat=newest}

\usepackage{calc}
\usepackage[nomessages]{fp}

\usepackage{datetime}
% Added by lyx2lyx
\setlength{\parskip}{\smallskipamount}
\setlength{\parindent}{0pt}

\makeatother

\usepackage{babel}
\providecommand{\exercisename}{Exercise}
\providecommand{\theoremname}{Theorem}

\begin{document}
\renewcommand{\author}{Dr.\,James Ashe}

\newcommand{\classNum}{439}

\newcommand{\classSec}{A}

\newcommand{\nameType}{solutions}

\newcommand{\examType}{Homework 4}

\newcommand{\Date}{\formatdate{6}{2}{2014}} %%\AdvanceDate[12] \today

%\newdate{my_date}{\day}{\month}{\year}%   
%\AdvanceDate[-5] 
%\displaydate{my_date}

%%First page's header%%%%%%%%%%%%%%%%%%%%%%
\fancypagestyle{firststyle} %%header settings for first pagr
{
	\fancyhead[L]{MTH \classNum{}(\classSec{}) \author{}\\
	\textbf{\examType{}} \Date{}}
	\fancyhead[C]{\textbf{\nameType}}
	\setlength{\headsep}{14pt}
}
%%%%%%%%%%%%%%%%%%%%%%%%%%%%%%%%%%%%%%%%%%

%%Next page's header%%%%%%%%%%%%%%%%%%%%%%%
%\setlist{nolistsep} %eliminates space between list items
\lhead{MTH \classNum{}(\classSec{})} 
\chead{\textbf{\nameType}} 
\cfoot{\thepage{}~of~\pageref{LastPage}} 
\setlength{\headsep}{5pt} 
%%%%%%%%%%%%%%%%%%%%%%%%%%%%%%%%%%%%%%%%%%%

%%Various settings; see the preamble for other settings%%%%%
%%\setenumerate[0]{leftmargin=12pt,itemindent=0pt,labelwidth=0pt}
\renewcommand{\labelenumi}{\textbf{\arabic{enumi}.}}
\renewcommand{\labelenumii}{\textbf{\alph{enumii}.}}
\thispagestyle{firststyle}
%%%%%%%%%%%%%%%%%%%%%%%%%%%%%%%%%%%%%%%%%%

\null

\vspace{1cm}

\setcounter{section}{3}

\Large{\textbf{Chapter 3}} \large

\textbf{\textbf{\examType} due: \formatdate{25}{3}{2014}}

\setcounter{thm}{13}

All problems \uline{without }a $\bigstar$ from section 3.2 and 3.3;
exercises 3.14-3.19. 
\begin{xca}
Find the distinct ideals of $\mathbb{Z}_{12}$. 

\begin{proof}
All of the ideals of $\mathbb{Z}_{12}$ are principal. It is easy
to directly check that the distinct ideals of $\mathbb{Z}_{12}$ are:
\begin{eqnarray*}
\left\langle 0\right\rangle  & = & \left\{ 0\right\} \\
\left\langle 1\right\rangle  & = & \left\{ 1,2,3,4,\dots,11,0\right\} \\
\left\langle 2\right\rangle  & = & \left\{ 2,4,6,8,10,0\right\} \\
\left\langle 3\right\rangle  & = & \left\{ 3,6,9,0\right\} \\
\left\langle 4\right\rangle  & = & \left\{ 4,8,0\right\} \\
\left\langle 6\right\rangle  & = & \left\{ 6,0\right\} 
\end{eqnarray*}

In general, the distinct ideals of $\mathbb{Z}_{n}$ will be all the
divisors of $n$ and $\left\langle 0\right\rangle $. To see this
suppose that $\gcd(n,m)=1$. Then $\left\langle m\right\rangle =\left\{ m,2m,3m,\dots km\right\} $.
Since $n$ and $m$ are relatively prime, $m(k-1)\not\equiv0\bmod(nm)$,
and $k=n$ since $\left|\left\langle m\right\rangle \right|\leq n$.
So $\gcd(m,n)=1$ implies that $\left\langle m\right\rangle =\left\langle 1\right\rangle $. 

Suppose $\mbox{lcm}(m,n)=k$. Then $\left\langle m\right\rangle =\left\{ m,2m,\dots,km\right\} $
where $m(k-1)\not\equiv0\bmod n$ since then $n|k-1$ and $k-1$ would
be the lowest common multiple. Also $n\nmid(am-bm)$ for any distinct
$0\leq a,b\leq k$ since then $m(a-b)$ would be the lowest common
multiple not $k$. So $\left\langle m\right\rangle $ has $k$ distinct
elements, and it follows that the distinct ideals of $\mathbb{Z}_{n}$
are the distinct divisors of $n$ and $0$. 
\end{proof}
\end{xca}

\begin{xca}
$\bigstar$%3.15
\end{xca}

\begin{xca}
$\bigstar$%3.16
\end{xca}

\begin{xca}
%3.17Let $R$ be a ring with zero element $0_{r}$. Determine the
elements of $R/0_{R}$ and $R/R$\@.

\begin{proof}
$(R/0_{R}):$ Let $a+0_{R}\in R/0_{R}$. For any $r\in R$, $a\equiv r\bmod0_{R}$
iff $a-r\in0_{R}\iff a=r$. So then $a+0_{R}=a$ for any $a\in R$.
Recall that $R/0_{R}=R$.

$(R/R):$ Let $a+R\in R/R$. For any $r\in R$, $a\equiv r\bmod R$
iff $a-r\in R$. $R$ is a ring so then $a-r\in R$ for all $a$ and
$r\in R$. This is true for $0_{R}$, i.e., $a-0_{R}\in R$ implying
that $a+R\equiv0+R$ for all $a\in R$. Thus $R/R$ has only $1$
coset, $0+R=R$. Recall that $R/R=0$. 
\end{proof}
\end{xca}

\newpage
\begin{xca}
%3.18Let $D$ be the ring of rational numbers (in reduced terms)
with odd denominators, and let $E$ be the set of elements in $D$
with even numerators. 

\begin{enumerate}
\item Prove that $E$ is an ideal in $D$.
\item Prove that $D/E$ has exactly two distinct cosets.
\end{enumerate}
\begin{proof}
Part 1: Note that any two rational numbers can be written so that
they have the same denominators. Let $\frac{2x}{2n+1},\frac{2y}{2n+1}\in E$
and $\frac{r}{2n+1}\in D$ for integers $x,y,n,r\in\mathbb{Z}$. $\frac{2x}{2n+1}-\frac{2y}{2n+1}=\frac{2(x-y)}{2n+1}$
and $\frac{2x}{2n+1}\frac{r}{2n+1}=\frac{2xr}{4n^{2}+4n+1}=\frac{2\left(xr\right)}{2(2n^{2}+2n)+1}\in E$.
So $E$ is an ideal in $D$ by the ideal test.

Part 2: Let $\frac{x}{2n+1}\in D$ then $x$ is either even or odd.
If it is even then $\frac{x}{2n+1}+E=0+E=E$ since $\frac{x}{2n+1}\in E$.
If it is odd then we can write $\frac{2m+1}{2n+1}$ for some $m\in\mathbb{Z}$.
Now let $\frac{2m'+1}{2n+1}$ for some $m'\in\mathbb{Z}$. $\frac{2m+1}{2n+1}-\frac{2m'+1}{2n+1}=\frac{2(m-m')}{2n+1}$
which has an even denominator thus $\frac{2m+1}{2n+1}+E=\frac{2m'+1}{2n+1}+E$.
So there are exactly 2 distinct cosets. 
\end{proof}
\end{xca}

\begin{itemize}
\item Note that since $1=\frac{1}{1}$, $\frac{2m+1}{2n+1}+E=1+E$. So the
two cosets can written as $0+E$ and $1+E$.
\end{itemize}

\begin{xca}
%3.19Let $C$ be the set of all polynomials with zero constant term
in $\mathbb{Z}[x]$.

\begin{enumerate}
\item Show that $C$ is the principal ideal $\left\langle x\right\rangle $. 
\item Show that $\mathbb{Z}[x]/C$ consists of an infinite number of distinct
cosets, one for each $n\in\mathbb{Z}$.
\end{enumerate}
\begin{proof}
Part 1: If $g(x)\in C$ then 
\[
g(x)=a_{n}x^{n}+a_{n-1}x^{n}+\cdots+a_{1}x=(x)(a_{n}x^{n-1}+a_{n-1}x^{n}+\cdots+a_{1})=x\cdot f(x)
\]
for some $f(x)\in\mathbb{Z}[x]$.

Part 2: Let $f_{1}(x)+C,f_{2}(x)+C\in\mathbb{Z}[x]/C$ for $f_{1}(x),f_{2}(x)\in\mathbb{Z}[x]$.
From part one, $f_{1}(x)+C=f_{1}(x)+\left\langle x\right\rangle $
and $f_{2}(x)+C=f_{2}(x)+\left\langle x\right\rangle $. If $f_{1}(x)-f_{2}(x)\in\left\langle x\right\rangle $
then $f_{1}(x)-f_{2}(x)$ must have a zero constant term ($x\cdot g(x)=f_{1}(x)-f_{2}(x)$
for some $g(x)\in\mathbb{Z}[x]$). 
\begin{eqnarray*}
f_{1}(x)-f_{2}(x) & = & \left[a_{n}x^{n}+a_{n-1}x^{n}+\cdots+a_{1}x+a_{0}\right]-\left[b_{n}x^{n}+b_{n-1}x^{n}+\cdots+b_{1}x+b_{0}\right]\\
 & = & (a_{n}-b_{n})x^{n}+(a_{n-1}-b_{n-1})x^{n}+\cdots+(a_{1}-b_{1})x+(a_{0}-b_{0})\\
 & = & (a_{n}-b_{n})x^{n}+(a_{n-1}-b_{n-1})x^{n}+\cdots+(a_{1}-b_{1})x+0
\end{eqnarray*}
Thus $a_{0}-b_{0}=0\Rightarrow a_{0}=b_{0}$. So $f_{1}(x)$ and $f_{2}(x)$
have the same constant term. So two cosets $f_{1}(x)+\left\langle x\right\rangle $
and $f_{2}(x)+\left\langle x\right\rangle $ are equivalent if and
only if $f_{1}(x)$ and $f_{2}(x)$ have the same constant terms.
These are polynomials with integer coefficients so there is a distinct
coset for each integer. 
\end{proof}
\end{xca}


\end{document}
